% Part: first-order-logic
% Chapter: sequent-calculus
% Section: derivations

\documentclass[../../../include/open-logic-section]{subfiles}

\begin{document}

\iftag{FOL}
      {\olfileid{fol}{seq}{der}}
      {\olfileid{pl}{seq}{der}}

\olsection{\usetoken{P}{derivation}}

\begin{explain}
We weten nu hoe een beginsequent eruitziet en welke afleidingsregels er
zijn. Vanuit die onderdelen bouwen we de !!^{derivation}s in de
sequentenrekening stap voor stap op. Dat heet inductief opbouwen. Een
!!{derivation} bestaat óf alleen uit een beginsequent, óf uit één of
twee !!{derivation}s met daarna één afleidingsstap.
\end{explain}

\begin{defn}[$\Log{LK}$ !!{derivation}]
Een \emph{$\Log{LK}$-!!{derivation}} van een sequent~$S$ is een boom
met eindig veel sequenten. Die boom moet aan de volgende voorwaarden
voldoen:
\begin{enumerate}
\item Bovenaan in de boom staan beginsequenten.
\item De onderste sequent van de boom is~$S$.
\item Elke sequent in de boom behalve $S$ is een uitgangspunt voor een
  correct toegepaste afleidingsregel. De conclusie van die regel staat
  direct onder die sequent in de boom.
\end{enumerate}
We noemen $S$ dan de \emph{eindsequent} van de !!{derivation}. Ook
zeggen we dat $S$ \emph{!!{derivable} is in $\Log{LK}$} (of
$\Log{LK}$-!!{derivable}).
\end{defn}

\begin{ex}
Elke beginsequent is zelf al !!a{derivation}. Dat geldt bijvoorbeeld
voor $!C \Sequent !C$. We krijgen daaruit een nieuwe !!{derivation} als
we bijvoorbeeld de regel $\LeftR{\Weakening}$ toepassen:
\begin{prooftree}
\Axiom$ \Gamma \fCenter \Delta $
\RightLabel{\LeftR{\Weakening}}
\UnaryInf$ !A, \Gamma \fCenter \Delta$
\end{prooftree}
Deze regel is algemeen: voor $!A$ mogen we elke !!{sentence} invullen,
dus ook~$!D$. We willen de premisse laten samenvallen met onze
beginsequent $!C \Sequent !C$. Daarom vullen we voor zowel $\Gamma$ als
$\Delta$ precies~$!C$ in. De conclusie wordt dan $!D, !C \Sequent !C$.
Het volgende is dus !!a{derivation}:
\begin{prooftree}
\Axiom$ !C \fCenter !C $
\RightLabel{\LeftR{\Weakening}}
\UnaryInf$ !D, !C \fCenter !C$
\end{prooftree}
Nu kunnen we nog een regel toepassen. Met $\LeftR{\Exchange}$ mogen we
twee !!{sentence}s aan de linkerkant van plaats laten wisselen. Het
volgende is dus ook een correcte !!{derivation}:
\begin{prooftree}
\Axiom$ !C \fCenter !C $
\RightLabel{\LeftR{\Weakening}}
\UnaryInf$ !D, !C \fCenter !C$
\RightLabel{\LeftR{\Exchange}}
\UnaryInf$ !C, !D \fCenter !C$
\end{prooftree}
De gebruikte regel heeft de volgende algemene vorm:
\begin{prooftree}
\Axiom$ \Gamma, !A, !B, \Pi \fCenter \Delta $
\RightLabel{\LeftR{\Exchange}}
\UnaryInf$ \Gamma, !B, !A, \Pi \fCenter \Delta,$
\end{prooftree}
In onze toepassing zijn $\Gamma$ en $\Pi$ allebei leeg en is $\Delta$
gelijk aan $!C$. Voor $!A$ en $!B$ vullen we respectievelijk $!D$ en
$!C$ in. Vrijwel op dezelfde manier zien we dat
\begin{prooftree}
\Axiom$ !D \fCenter !D $
\RightLabel{\LeftR{\Weakening}}
\UnaryInf$ !C, !D \fCenter !D$
\end{prooftree}
!!a{derivation} is. We kunnen deze twee !!{derivation}s nu combineren
met $\RightR{\land}$. Die regel heeft de vorm
\begin{prooftree}
\Axiom$\Gamma \fCenter \Delta, !A$
\Axiom$ \Gamma \fCenter \Delta, !B$
\RightLabel{\RightR{\land}}
\BinaryInf$ \Gamma \fCenter \Delta, !A \land !B $
\end{prooftree}
De premissen van deze regel moeten precies gelijk zijn aan de laatste
sequenten van de twee !!{derivation}s. Daarom is $\Gamma$ hier $!C,
!D$, is $\Delta$ leeg, is $!A$ gelijk aan $!C$ en is $!B$ gelijk aan
$!D$. Alleen dan klopt de afleidingsstap. De conclusie is dus $!C, !D
\Sequent !C \land !D$.
\begin{prooftree}
\Axiom$ !C \fCenter !C $
\RightLabel{\LeftR{\Weakening}}
\UnaryInf$ !D, !C \fCenter !C$
\RightLabel{\LeftR{\Exchange}}
\UnaryInf$ !C, !D \fCenter !C$
\Axiom$ !D \fCenter !D $
\RightLabel{\LeftR{\Weakening}}
\UnaryInf$ !C, !D \fCenter !D$
\RightLabel{\RightR{\land}}
\BinaryInf$ !C, !D \fCenter !C \land !D $
\end{prooftree}
We kunnen de twee premissen ook omdraaien. Dan is $!A$ gelijk aan $!D$
en $!B$ aan~$!C$.
\begin{prooftree}
\Axiom$ !D \fCenter !D $
\RightLabel{\LeftR{\Weakening}}
\UnaryInf$ !C, !D \fCenter !D$
\Axiom$ !C \fCenter !C $
\RightLabel{\LeftR{\Weakening}}
\UnaryInf$ !D, !C \fCenter !C$
\RightLabel{\LeftR{\Exchange}}
\UnaryInf$ !C, !D \fCenter !C$
\RightLabel{\RightR{\land}}
\BinaryInf$ !C, !D \fCenter !D \land !C $
\end{prooftree}
\end{ex}
\end{document}
